% The exit-code contract. A stage is a subprocess, so its exit code is the
% only thing the orchestrator can route on -- and code 3 is overloaded,
% which is why it branches again on the stage that produced it.
%
% Column spacing must exceed the widest box or the columns collide; the
% x unit below is set from the box width, not guessed.
\begin{tikzpicture}[font=\plexsans\scriptsize, x=4.45cm, y=-1.16cm]

  \tikzset{
    ec/.style={box, text width=3.15cm, minimum height=1.00cm},
    dec/.style={box, fill=mist, draw=slate, text width=1.35cm,
                minimum height=0.86cm},
    term/.style={box, rounded corners=8pt, text width=1.35cm,
                 minimum height=0.72cm, font=\plexsans\bfseries\scriptsize}}

  \node[dec] (c) at (0,2) {exit code};

  \node[ec, pass] (ok)    at (1,0) {\textbf{0}\enspace Work done,\\nothing to flag};
  \node[dec]      (w)     at (1,1.6) {\textbf{3}\\which stage?};
  \node[ec, stop] (miss)  at (1,3.1) {\textbf{2}\enspace Required upstream\\artifact missing};
  \node[ec, stop] (crash) at (1,4.3) {\textbf{1}\enspace Unhandled runtime\\or config error};

  \node[ec, gate] (qual) at (2,0.75) {\code{07 · 08 · 09 · 12}\\[1pt]
        \textbf{Data quality}\\{\plexsanslight output written, needs review}};
  \node[ec, stop] (part) at (2,2.55) {\code{03}\\[1pt]
        \textbf{Extraction incomplete}\\{\plexsanslight partials kept for resume}};

  \node[term, pass] (cont) at (3,0.0) {continue};
  \node[term, stop] (halt) at (3,3.6) {stop};

  \foreach \t in {ok, w, miss, crash} \draw[flow] (c) -- (\t);
  \draw[flow] (w) -- (qual);
  \draw[flow] (w) -- (part);
  \draw[flow] (ok)   -- (cont);
  \draw[flow] (qual) -- (cont);
  \foreach \s in {part, miss, crash} \draw[flow] (\s) -- (halt);

\end{tikzpicture}
